% 按内容保留有价值的字段；没有实质内容的复盘字段可以删除。
\subsection{可检索的题型名称}

\problemAnchor{MATH1-CALC-0000}
\begin{problemBox}
\begin{problemMeta}
\textbf{题目编号：} \problemRef{MATH1-CALC-0000} \quad
\textbf{索引：} \problemIndexRef{MATH1-CALC-0000} \quad
\textbf{来源：} 用户提供 / 未注明来源

\textbf{题型：} 题型名称 \quad
\textbf{难度：} 基础 / 中等 / 较难 / 压轴

\textbf{知识点：} \knowledgeRef[knowledge-key]{知识点名称}
\end{problemMeta}

\textbf{题目：}

在此写题面。影响求解的定义域、参数限制或 OCR 校注写在题面附近。
\end{problemBox}

\begin{solutionBox}
\textbf{思路入口：}

写出识别信号和自然入口。

\textbf{详细解答：}

展开决定正确性的推导，并检查本题实际使用的条件。

\textbf{答案：}

写出最终结论。
\end{solutionBox}

\knowledgeAnchor[knowledge-key]{知识点名称}
\begin{knowledgeBox}
\textbf{识别信号：} 说明何时想到本方法。

\textbf{核心方法：} 提炼可迁移动作和适用条件。

\textbf{考场写法：} 保留真实拿分步骤。

\textbf{变式与迁移：} 只记录有复习价值的母题或变式。
\end{knowledgeBox}

\begin{mistakeBox}
\textbf{易错点：} 说明错因、后果和避免方法。

\textbf{复盘提醒：} 给出一个可执行的复盘动作。
\end{mistakeBox}
